\section{Résultats de rendement}
\subsection{Production finale par plot}
\begin{figure}[H]
    \centering
    \imgstd{images/z15_navet_temoin_chap5.jpg}
    \caption{Rendement obtenu sur le NAV T. }
    \label{fig:rendement-navT}
\end{figure}
Sur le témoin, nous n’avons que trois plants résistants, qui ont fini par produire malgré leur retard de cycle. Cependant, les tubercules obtenus ne sont pas sains.

\begin{figure}[H]
    \centering
    \imgstd{images/z16_navet_test_chap5.jpg}
    \caption{Rendement obtenu sur le NAV TS. }
    \label{fig:rendement-navTS}
\end{figure} 
Sur le plot test, la récolte a été satisfaisante : les \textbf{tubercules étaient  sains et bien développés} (en taille et en diamètre), avec un \textbf{rendement de 4{,}8~kg} au total. C’est la première fois qu’un tel résultat est obtenu en plein champ, hors serre. \textbf{Depuis la mise en place du champ d'expérimentation Agro TIC} par nos aînés de la première ou deuxième promotion, \textbf{aucune culture de navet n’avait donné un tel résultat en plein champ}. Ce résultat soutient l’intérêt de l’approche testée.

\begin{figure}[H]
    \centering
    \imgstd{images/z17_relation_longueur_Diametre.jpeg}
    \caption{Relation entre la longueur et le diamètre de quelques tubercules de navet mesurés réalisés après récolte. }
    \label{fig:long-Diam-navet}
\end{figure} 
La croissance des tubercules sur la partie témoin est plus variable, probablement en raison de la salinité, de la compacité du sol, des attaques de parasites et du stress lié à l'eau. La production est plus homogène dans la partie test.

\begin{figure}[H]
    \centering
    \imgstd{images/z18_concurence.jpeg}
    \caption{Relation de comparaison entre quelques tubercules de navet présentant une concurrence. }
    \label{fig:navet-concurence}
\end{figure} 
La compétition entre plants est déséquilibré : la plante « Max » bénéficie d’un accès supérieur 
aux ressources au détriment de la plante « Min », les différences pouvant aller jusqu'à 50~\% en longueur tubercule et diamètre tubercule. Ceci implique une plantation trop dense.

\begin{figure}[H]
    \centering
    \imgstd{images/z19_rendement_laitue_test.jpg}
    \caption{Rendement obtenu sur la LAI TS. }
    \label{fig:rendement-laitue}
\end{figure} 
Pour la laitue, il n'y a que le test qui a donné une récolte. Les témoins n’ont pas survécu.

\begin{figure}[H]
    \centering
    \imgstd{images/z20_rendement_gombo_test.jpg}
    \caption{Rendement obtenu sur le GOM TS par intervalle d’environ 5 jours. }
    \label{fig:rendement-gombo}
\end{figure} 
Pour le gombo test la récolte est continue, pas une seule fois mais plutôt tous les 4 à 7 jours pour un nombre variant entre 7 à 10 fruits.

\subsection{Comparaison des rendements entre T et TS}

\begin{table}[H]
    \centering
    \caption{Comparaison des rendements entre plot T et plot TS}
    \label{tab:comparaison_rendement}
    \begin{tabular}{|p{2.4cm}|p{5.6cm}|p{6cm}|}
        \hline
        \textbf{Spéculation} & \textbf{Plot T} & \textbf{Plot TS} \\
        \hline
        Gombo & Aucune récolte : un plant seul a survécu sans fructification & Récolte régulière après 2 mois, presque 3 mois ; les plants survivants produisent environ 7 à 10 fruits par période de 5 à 7 jours maximum \\
        \hline
        Navet & 3 tubercules récoltés hors cycle (plus de 2 mois) avec une peau anormale & Récolte effectuée dans le cycle normal, avec un rendement de 4,8~kg sur 2~m\textsuperscript{2} \\
        \hline
        Laitue & Aucune récolte : mort totale de tous les plants & Malgré les pertes par fonte, récolte de plus d'une dizaine de pommes \\
        \hline
    \end{tabular}
\end{table}

L’écart de rendement entre plots T et TS est cohérent avec les observations précédentes : sur les plots T, la dégradation de la structure et la salinité du sol ont probablement affecté dès le départ la levée et l’installation des racines, ce qui s’accorde avec la mortalité de la laitue, le retard et la faible qualité du navet, et l’absence quasi totale de production du gombo. Sur les plots TS, le sol amendé et l’irrigation pilotée ont permis un développement jusqu’à une production effective, malgré les pertes liées aux maladies fongiques. L’amendement semble ainsi avoir conditionné non seulement l’état du sol, mais aussi la capacité des cultures à survivre et à produire.







